\chapter{动量与碰撞}

动量方法最擅长处理短时强相互作用、系统内力复杂但外力冲量可控的问题。对于变质量系统，动量守恒的使用更需要谨慎，必须分清系统边界与质量通量。本章从常见碰撞与反冲出发，进一步延伸到火箭、链条和漏沙系统。

\section*{经典题}

\begin{classic}
质量分别为 $m_1,m_2$ 的两小球在光滑直线上对心碰撞。碰前速度分别为 $u_1,u_2$，碰后速度为 $v_1,v_2$。设碰撞为完全弹性碰撞。

(1) 写出动量守恒与机械能守恒方程；

(2) 解出碰后速度；

(3) 讨论 $m_1=m_2$ 的特殊情形。

\begin{solution}
完全弹性碰撞满足
\[
m_1u_1+m_2u_2=m_1v_1+m_2v_2,
\]
\[
\frac12 m_1u_1^2+\frac12 m_2u_2^2=\frac12 m_1v_1^2+\frac12 m_2v_2^2.
\]
联立可得经典结果
\[
v_1=\frac{m_1-m_2}{m_1+m_2}u_1+\frac{2m_2}{m_1+m_2}u_2,
\]
\[
v_2=\frac{2m_1}{m_1+m_2}u_1+\frac{m_2-m_1}{m_1+m_2}u_2.
\]

也可利用弹性碰撞相对速度反号关系
\[
v_1-v_2=-(u_1-u_2)
\]
与动量守恒联立快速求解。

若 $m_1=m_2$，则上式化为
\[
v_1=u_2,
\qquad
v_2=u_1.
\]
即两球交换速度。这正是台球中最常见的“同质量对心弹碰”规律。
\end{solution}
\end{classic}

\begin{classic}
一门炮放在光滑水平轨道上，炮身质量为 $M$，炮弹质量为 $m$。炮弹以相对炮口速度 $u$ 水平射出。求炮身反冲速度与炮弹对地速度。

\begin{solution}
设炮弹向右、炮身向左。以地面为参考系，发射前系统静止，总动量为零。发射后若炮身速度为 $V$（向左取负），则炮弹对地速度为
\[
v=u+V.
\]
因为“相对炮口速度为 $u$”意味着炮弹相对炮身的速度差为 $u$。

由动量守恒
\[
0=mv+MV=m(u+V)+MV.
\]
解得
\[
V=-\frac{mu}{M+m}.
\]
负号表示向左反冲。炮弹对地速度为
\[
v=u+V=u-\frac{mu}{M+m}=\frac{M}{M+m}u.
\]

若 $M\gg m$，则
\[
\abs{V}\ll u,
\qquad v\approx u,
\]
即炮身后退很慢，炮弹速度几乎等于炮口给出的相对速度。
\end{solution}
\end{classic}

\begin{classic}
一静止炸弹在空中炸裂成质量分别为 $m_1,m_2$ 的两块碎片。若其中一块沿 $x$ 方向速度为 $\vb v_1=v_1\vu x$，求另一块速度 $\vb v_2$。并说明如何用冲量观点理解这一结果。

\begin{solution}
爆炸过程极短，外力冲量（如重力冲量）相对内爆炸力产生的冲量可忽略，因此系统总动量守恒。爆炸前静止，总动量为零：
\[
m_1\vb v_1+m_2\vb v_2=\vb 0.
\]
故
\[
\vb v_2=-\frac{m_1}{m_2}\vb v_1=-\frac{m_1v_1}{m_2}\vu x.
\]

若从微积分角度看，设爆炸内力非常大但持续时间很短，系统总外力冲量为
\[
\Delta \vb p=\int_{t_1}^{t_2}\vb F_\text{ext}\,\dd t\approx \vb 0.
\]
因此总动量在爆炸前后近似不变。这个结论与内力如何复杂分布无关，只取决于外力总冲量是否可忽略。
\end{solution}
\end{classic}

\section*{竞赛题}

\begin{competition}
一条均匀柔软绳，总长 $L$、总质量 $M$，初始时有长度 $x$ 垂在光滑桌边下方，其余部分在桌面上静止。设绳开始在重力作用下滑落，求 $x$ 的运动方程，并求绳端刚开始运动时的加速度。

\begin{solution}
取整条绳为研究对象，线密度
\[
\lambda=\frac{M}{L}.
\]
设所有绳元速度大小相同为 $v$，桌面光滑。系统沿绳方向的驱动力仅为悬垂部分重力
\[
F=\lambda x g.
\]
总质量为 $M$，故由牛顿第二定律
\[
M\dv{v}{t}=\lambda x g.
\]
又因为 $v=\dot x$，所以
\[
M\ddot x=\lambda x g.
\]
代入 $\lambda=M/L$，得
\[
\ddot x=\frac{g}{L}x.
\]
这是一条线性常微分方程，其解为指数/双曲函数形式：
\[
x(t)=C_1\mathrm e^{\sqrt{g/L}t}+C_2\mathrm e^{-\sqrt{g/L}t}.
\]

若初始悬垂长度为 $x_0$，且初速为零，则
\[
x(0)=x_0,
\qquad \dot x(0)=0
\]
给出
\[
x(t)=x_0\cosh\qty(\sqrt{\frac{g}{L}}t).
\]
于是刚开始运动时的加速度为
\[
a_0=\ddot x(0)=\frac{g}{L}x_0.
\]
可见垂下越长，初始驱动力越大，系统越容易被“拖”下桌边。
\end{solution}
\end{competition}

\begin{competition}
一条均匀链条以速度 $v$ 竖直落到地面，单位长度质量为 $\lambda$。求地面对链条的支持力，设地面上已堆积部分静止，忽略反弹。

\begin{solution}
设在某时刻已有长度 $x$ 落在地面上，则静止堆积部分重量为
\[
\lambda x g.
\]
此外，连续有新链条以速度 $v$ 撞地并在极短时间内减速为零，因此地面还要提供动量流对应的额外力。

单位时间内落地的质量为
\[
\dv{m}{t}=\lambda v.
\]
这些质量的动量变化率大小为
\[
\dv{p}{t}=\lambda v\cdot v=\lambda v^2.
\]
故地面对链条的支持力为
\[
N=\lambda x g+\lambda v^2.
\]
第一项用于托住已经堆在地上的链条，第二项用于把新到达的链条元从速度 $v$ 瞬间降到 $0$。很多同学只算重量项而漏掉动量流项，这是变质量/连续介质问题最常见的错误。
\end{solution}
\end{competition}

\begin{competition}
质量为 $m$ 的子弹以速度 $u$ 射入静止在光滑水平面上的木块，木块质量为 $M$。若子弹穿出后速度为 $v$，求：

(1) 木块获得的速度；

(2) 系统损失的机械能；

(3) 当子弹恰嵌入木块时结果如何变化。

\begin{solution}
设木块获得速度为 $V$。水平方向动量守恒：
\[
mu=mv+MV.
\]
故
\[
V=\frac{m(u-v)}{M}.
\]

初始机械能
\[
E_i=\frac12 mu^2.
\]
末态机械能
\[
E_f=\frac12 mv^2+\frac12 MV^2
=\frac12 mv^2+\frac{m^2(u-v)^2}{2M}.
\]
故机械能损失
\[
\Delta E=E_i-E_f
=\frac12 m(u^2-v^2)-\frac{m^2(u-v)^2}{2M}.
\]
也可整理为
\[
\Delta E=\frac{m(u-v)}{2}\qty[u+v-\frac{m(u-v)}{M}].
\]

若子弹恰嵌入木块，则 $v=V$。由动量守恒
\[
mu=(M+m)V,
\qquad
V=\frac{m}{M+m}u.
\]
机械能损失为
\[
\Delta E=\frac12 mu^2-\frac12(M+m)V^2
=\frac12\frac{mM}{M+m}u^2.
\]
这说明完全非弹性碰撞的能量损失最大。
\end{solution}
\end{competition}

\section*{大学物理题}

\begin{university}
推导齐奥尔科夫斯基火箭方程。设火箭瞬时质量为 $m$，速度为 $v$，相对火箭喷出燃气速度大小为 $u$，忽略外力。

\begin{solution}
考虑极短时间 $\dd t$ 内，火箭喷出微元质量 $-\dd m$（其中 $\dd m<0$）。在时刻 $t$，火箭动量为
\[
p_i=mv.
\]
到 $t+\dd t$，火箭质量变为 $m+\dd m$，速度变为 $v+\dd v$；喷出的燃气质量为 $-\dd m$，对地速度为 $v-u$。于是末态总动量为
\[
p_f=(m+\dd m)(v+\dd v)+(-\dd m)(v-u).
\]
忽略二阶小量 $\dd m\dd v$，展开得
\[
p_f=mv+m\dd v+v\dd m-v\dd m+u\dd m=mv+m\dd v+u\dd m.
\]
外力忽略时动量守恒，$p_f=p_i$，故
\[
m\dd v+u\dd m=0.
\]
整理为
\[
\dd v=-u\frac{\dd m}{m}.
\]
从初始质量 $m_0$、初始速度 $v_0$ 积分到末态 $m,v$，得
\[
\int_{v_0}^{v}\dd v'=-u\int_{m_0}^{m}\frac{\dd m'}{m'}.
\]
因此
\[
v-v_0=u\ln\frac{m_0}{m}.
\]
这就是齐奥尔科夫斯基火箭方程。

若考虑重力和阻力，只需在微分形式中加入外力项：
\[
m\dv{v}{t}=-u\dv{m}{t}+F_\text{ext}.
\]
但火箭方程的核心对数结构仍来自“喷气速度恒定、质量持续减少”的积分效应。
\end{solution}
\end{university}

\begin{university}
用质心系分析一维碰撞。说明为何在质心系中，完全弹性碰撞前后速度只会反向而大小不变。

\begin{solution}
设两粒子质量为 $m_1,m_2$，实验室系速度为 $u_1,u_2$。质心速度为
\[
V_C=\frac{m_1u_1+m_2u_2}{m_1+m_2}.
\]
在质心系中，碰前速度为
\[
u_1'=u_1-V_C,
\qquad
u_2'=u_2-V_C.
\]
按定义有
\[
m_1u_1'+m_2u_2'=0,
\]
所以两者动量等大反向。

若碰撞为完全弹性，则质心系中总动量守恒、总动能也守恒。既然总动量始终为零，碰后仍需满足
\[
m_1v_1'+m_2v_2'=0.
\]
又因总动能守恒，速度大小只能与碰前相同，否则无法同时满足两条守恒律。于是必有
\[
v_1'=-u_1',
\qquad
v_2'=-u_2'.
\]
也就是说，在质心系中弹性碰撞的效果极其简单：每个粒子速度方向反转，大小不变。

再加回质心速度即可回到实验室系：
\[
v_1=v_1'+V_C,
\qquad
v_2=v_2'+V_C.
\]
这就是许多碰撞公式最简洁的推导路径。
\end{solution}
\end{university}

\begin{university}
一质量为 $m_1$ 的小球以速度 $u$ 沿 $x$ 轴运动，与静止的质量为 $m_2$ 的小球发生二维完全弹性碰撞。碰后 $m_1$ 速度大小为 $v_1$，偏转角为 $\theta$，$m_2$ 速度大小为 $v_2$，偏转角为 $\phi$。

(1) 写出动量守恒方程；

(2) 写出能量守恒方程；

(3) 当 $m_1=m_2$ 时证明两球碰后速度方向互相垂直。

\begin{solution}
二维动量守恒分量式为
\[
m_1u=m_1v_1\cos\theta+m_2v_2\cos\phi,
\]
\[
0=m_1v_1\sin\theta-m_2v_2\sin\phi.
\]
机械能守恒为
\[
\frac12 m_1u^2=\frac12 m_1v_1^2+\frac12 m_2v_2^2.
\]

当 $m_1=m_2=m$ 时，方程化为
\[
u=v_1\cos\theta+v_2\cos\phi,
\qquad
0=v_1\sin\theta-v_2\sin\phi,
\qquad
u^2=v_1^2+v_2^2.
\]
把前两式平方后相加：
\[
u^2=v_1^2+v_2^2+2v_1v_2\cos(\theta+\phi).
\]
与能量守恒比较，可得
\[
2v_1v_2\cos(\theta+\phi)=0.
\]
非平凡碰撞下 $v_1,v_2\ne 0$，故
\[
\cos(\theta+\phi)=0
\quad\Rightarrow\quad
\theta+\phi=90^\circ.
\]
因此两球碰后速度方向互相垂直。这是同质量二维弹性碰撞的经典结论。
\end{solution}
\end{university}

\section*{普通/微分方程题}

\begin{problem}
一节漏沙车厢初始总质量为 $M_0$，其中沙子以恒定速率 $\mu$ 从车底垂直漏下。车厢在光滑水平轨道上，初速度为 $v_0$。求其运动方程并讨论速度是否改变。

\begin{solution}
关键在于：漏下的沙子在离开车厢前与车厢具有相同的水平速度，且漏沙方向竖直，因此没有额外水平相对喷射速度。取“车厢加当前仍在车上的沙”为系统，水平方向无外力。

在短时间 $\dd t$ 内，漏出质量 $\mu\dd t$。由于漏出的沙子离开瞬间水平速度仍为 $v$，故系统的水平动量并不会因为质量减少而产生额外推力项。微分形式可写成
\[
\dv{}{t}(Mv)=0.
\]
但这里必须谨慎：该式中的“逸出质量”带走的正是其本来就具有的动量，因此对剩余车厢系统应用雷诺输运思想后，可得
\[
M\dv{v}{t}=0.
\]
所以
\[
v(t)=v_0.
\]

质量随时间变化为
\[
M(t)=M_0-\mu t,
\qquad 0\le t<\frac{M_0}{\mu}.
\]
但速度始终不变，位移为
\[
x(t)=x_0+v_0 t.
\]
这与火箭问题形成鲜明对比：火箭喷气有相对速度，漏沙则无水平相对速度，因此不会产生推进或减速效应。
\end{solution}
\end{problem}

\begin{problem}
概述变质量系统的一般理论。设主体质量为 $m(t)$、速度为 $\vb v$，单位时间有质量流率 $\dot m$ 穿过系统边界，相对主体的流入/流出速度为 $\vb u_\text{rel}$。写出一般动量方程，并说明使用时最容易犯的错误。

\begin{solution}
对开放系统，不能直接把 $m\dv*{\vb v}{t}=\vb F_\text{ext}$ 生搬硬套，因为系统质量在变，且穿越边界的质量携带动量通量。更一般的动量平衡应写为
\[
\vb F_\text{ext}=\dv{}{t}(m\vb v)-\vb u_\text{rel}\dot m,
\]
其中符号约定要与 $\dot m$ 的正负统一。若取 $\dot m<0$ 表示流出，则常见火箭方程写成
\[
m\dv{\vb v}{t}=\vb F_\text{ext}+\vb u_\text{rel}(-\dot m).
\]

这个公式的核心是：穿越边界的质量相对于主体若有速度差，就会对主体产生反作用；若没有相对速度差，质量变化本身不一定改变主体速度。

最常见的错误有三类：
\begin{enumerate}
  \item 忽略质量通量项，错误地写成 $m\dot v=F$；
  \item 把相对速度误写成对地速度，导致符号和大小都错；
  \item 没有先明确“系统边界”，结果同一道题里混用了不同系统。
\end{enumerate}

解决变质量题的标准流程是：先选系统，再写边界质量流向与相对速度，最后列动量平衡。只要这三步清楚，火箭、链条、漏沙、雨滴等问题都可以放进统一框架中处理。
\end{solution}
\end{problem}
